\section{Introduction}
\label{sec:introduction}

\IEEEPARstart{T}{he} ancient Roman town of Herculaneum, buried by the catastrophic eruption of Mount Vesuvius in 79~AD, housed an extensive library of papyrus scrolls-the only intact collection of literary and philosophical texts surviving from classical antiquity~\cite{parsons1936}. These scrolls, carbonized by volcanic pyroclastic flows at temperatures exceeding 300\textdegree C, became brittle, blackened cylinders indistinguishable from charcoal. For over two centuries, attempts at physical unrolling resulted in their disintegration, destroying the very texts they were meant to reveal~\cite{janko2002herculaneum}. The challenge of reading these scrolls without physical contact has motivated decades of research at the intersection of imaging science, materials analysis, and computational methods.

The advent of high-resolution X-ray computed tomography (CT) offered a transformative approach: scanning the intact scrolls volumetrically and computationally ``unwrapping'' the layered papyrus surfaces to produce readable images~\cite{seales2005digital, seales2023educelab}. In this paradigm, 3D CT volumes are first segmented to identify individual papyrus layers, which are then flattened into 2D surface images. The critical final step-\emph{ink detection}-involves identifying regions where ancient carbon-based ink deposited on the papyrus surface creates subtle density variations in the CT signal~\cite{mocella2015revealing, liu2024ink}.

The Vesuvius Challenge, launched in March 2023 by Brent Seales and Nat Friedman, galvanized the global research community by offering substantial prizes for reading text from intact Herculaneum scrolls using CT data~\cite{vesuvius2023challenge}. The challenge produced remarkable results: in February 2024, the Grand Prize was awarded to a team (Youssef Nader, Luke Farritor, and Julian Schilliger) for successfully reading multiple columns of Greek text from Scroll~1 (PHerc.Paris.4)~\cite{nader2024gp, farritor2023first}. Their solution employed a TimeSformer-based architecture~\cite{bertasius2021timesformer} trained on carefully curated ink labels across 26 depth layers of the virtually unwrapped papyrus surface.

\subsection{The Competition Code Problem}

While the Vesuvius Challenge achieved a historic scientific breakthrough, its outputs suffer from a fundamental limitation: \emph{competition code is optimized to win, not to be understood}. The winning repositories\footnote{\url{https://github.com/younader/Vesuvius-Grandprize-Winner}} \footnote{\url{https://github.com/jaredlandau/Vesuvius-Grandprize-Winner-Plus}} contain functional implementations that produce correct results, but lack the documentation, theoretical justification, and systematic ablation studies expected of scientific contributions. Key design decisions-such as which depth layers to use, how many layers are sufficient, and why specific architectures were chosen-are embedded as hard-coded constants without explanation. Although the winning team published a blog post describing their approach~\cite{naderink2024blog}, no formal academic paper accompanied the code to document and justify the methodology.

This absence creates a critical gap: the research community has a proven ink detection system but no scientific baseline against which to measure improvements. Practitioners must reverse-engineer the code to understand its behavior, and any modifications risk degrading performance because the sensitivity of the system to its hyperparameters is unknown.

\subsection{The Z-Layer Selection Problem}

Among the many design choices in the ink detection pipeline, \emph{depth layer selection}-the choice of which and how many Z-layers from the CT volume to present to the model is particularly consequential yet poorly understood. In a typical pipeline, each papyrus segment is represented by a stack of 65 grayscale images, each corresponding to a different depth offset from the computationally estimated papyrus surface. These layers span approximately 514~\textmu m of depth at a voxel resolution of 7.91~\textmu m per layer.

The ink signal is expected to reside in a narrow band near the papyrus surface, where carbon-based ink particles create density variations detectable by X-ray CT~\cite{mocella2015revealing, bukreeva2016virtual}. However, the exact depth distribution of the ink signal depends on multiple factors: the physical thickness of the ink layer, the accuracy of the surface estimation algorithm, and the local geometry of the papyrus sheet (which may be warped, compressed, or damaged). Selecting too few layers risks missing the ink signal entirely; selecting too many introduces noise from irrelevant tissue layers and increases computational cost proportionally.

The Grand Prize solution uses layers 17-42 (26 layers spanning ${\sim}206$~\textmu m), an empirically chosen range that was never systematically justified. Whether this is optimal, sufficient, or merely adequate remains an open question. Furthermore, whether these settings generalize across different scroll segments-each with potentially different surface geometry and ink characteristics-has not been tested.

\subsection{Research Objectives}

This work addresses the gap between competition code and scientific understanding by conducting a rigorous, systematic study of Z-layer sampling for ink detection in Herculaneum papyri. Our research objectives are:

\begin{enumerate}
    \item \textbf{Formalize the pipeline:} Reverse-engineer the Grand Prize winning solution and document its methodology with sufficient detail for independent reproduction.
    
    \item \textbf{Quantify layer-count sensitivity:} Determine how ink detection performance scales with the number of input depth layers, identifying the point of diminishing returns.
    
    \item \textbf{Identify optimal depth windows:} Sweep the layer selection window across the papyrus cross-section to locate the most informative depth range.
    
    \item \textbf{Measure individual layer contributions:} Perform leave-one-out sensitivity analysis to identify which specific layers carry the strongest ink signal.
    
    \item \textbf{Assess generalization:} Evaluate whether a single Z-layer configuration works robustly across diverse scroll segments without manual tuning.
    
    \item \textbf{Characterize failure modes:} Systematically analyze cases where the default configuration fails, proposing mechanistic hypotheses for each failure.
\end{enumerate}


